\section{Methods}
\label{section:methods}

This chapter defines the pushing task and the five models that are compared on
it. It begins with the simulation environment and the formal Markov Decision
Process, including the object-relative push primitive that serves as the action.
It then explains why the hand-tuned improvement reward of an initial baseline
fails, and derives the potential-based reward that replaces it. The reward is
presented first as separate position and orientation potentials, then as a single
Special Euclidean group $\mathrm{SE}(2)$ distance. The chapter closes with the
five models and with an analysis of why potential-based shaping partially assists
self-play without closing the gap to single-agent training.

\subsection{Simulation Environment}
\label{sec:env}

All experiments are performed in NVIDIA Isaac Lab 2.3.0, built on Isaac Sim
5.1.0, using the PhysX physics engine on the Graphical Processing Unit (GPU)
\cite{Mittal2023IsaacLab, Makoviychuk2021IsaacGym}. The agent is a Universal
Robots UR5e arm fitted with a Robotiq 2F-140 parallel-jaw gripper. The
manipulated object is a T-shaped block resting on a tabletop, and the goal is
displayed as a coloured ghost of the same block at the target pose.
Figure~\ref{fig:task_schematic} shows the setup.

\begin{figure}[H]
    \centering
    \includegraphics[width=0.55\textwidth]{diagram_task_schematic.png}
    \caption{The planar-pushing task. A UR5e arm pushes a T-block (blue) across a
    tabletop toward a goal pose (orange). Goals are defined in $\mathrm{SE}(2)$
    and combine a target position and a target orientation.}
    \label{fig:task_schematic}
\end{figure}

The principal advantage of Isaac Lab for this study is GPU-batched parallelism.
Many copies of the environment are stepped at the same time on a single device,
which produces the large sample batches on-policy learning requires. The
single-agent and curriculum models are trained with 528 parallel environments,
which yields a batch of roughly 7{,}920 push transitions per Proximal Policy
Optimization (PPO) update. A hardware upper bound of one RTX Pro 6000 GPU
(96\,GB of memory) sets the compute budget for every model, so all methods are
compared on equal terms. Inverse Kinematics (IK) for the arm is solved with
cuRobo 0.7.5, a GPU-accelerated batched IK solver.

\subsection{The Planar-Pushing Markov Decision Process}
\label{sec:mdp}

The task is modelled as an episodic Markov Decision Process (MDP) with a
macro-action, so that a single decision corresponds to one complete push rather
than to a low-level joint command.

\paragraph{State.}
The state is a 28-dimensional vector composed of four groups. These are the
end-effector pose (6 values), the object pose and velocity (14 values), the goal
pose (6 values), and the two scalar position and orientation errors between object
and goal. As explained in Section~\ref{sec:pomdp}, the current observation does
not report the underlying contact state. The problem is therefore formally
partially observable, and every policy is recurrent.

\paragraph{Action.}
The policy does not command joint torques. It selects an \emph{object-relative
push primitive} described by four parameters: a radial approach offset $r$, an
approach angle $\phi$ in the object frame, a push length $\ell$, and a push
direction $\theta$ in the world frame. Each parameter is discretised into 21 bins,
so the action is a four-dimensional MultiCategorical variable with $21^4 =
194{,}481$ discrete values. Parameterising the approach relative to the object,
rather than in absolute table coordinates, guarantees contact wherever the object
lies, which raises the contact rate to roughly 95\,\%. Once decoded, the primitive
runs as a five-phase trajectory, approach, descend, push, retract, and return,
spanning 72 physics substeps. Each waypoint is mapped to joint targets by the
cuRobo IK solver.

\paragraph{Transition.}
The object moves according to the quasi-static pushing mechanics of
Section~\ref{sec:pushing}. The limit-surface normality condition of
Equation~\ref{eq:normality} couples translation and rotation: almost every push
changes both the position and the orientation of the block. This coupling is a
property of the environment, not a modelling choice. It is the reason the task
cannot be decomposed into independent position and orientation sub-problems.

\paragraph{Reward.}
The reward is the subject of Sections~\ref{sec:frac_reward} to \ref{sec:se2}. It
is the single component that differs across the models compared in this thesis,
which is what makes the comparison a controlled one.

\paragraph{Success criterion.}
A push episode is counted as successful when the object is within
$0.05\,\text{m}$ of the goal position \emph{and} within $0.2\,\text{rad}$ of the
goal orientation. This combined gate is applied identically in training and
evaluation. The two conditions must hold at the same time, and the physics couples
the two axes. The combined gate is therefore the central difficulty of the task,
and the quantity every model is ultimately measured against.

\subsection{Why the Improvement Reward Fails}
\label{sec:frac_reward}

An initial baseline used a hand-tuned reward based on fractional improvement. It
rewards each push in proportion to the fraction of the remaining error it removes.
Let $d_{\text{prev}}$ and $d_{\text{now}}$ be the position error before and after
a push, and $y_{\text{prev}}$, $y_{\text{now}}$ the orientation error. The reward
is then
\begin{equation}
    R = \alpha\,\frac{d_{\text{prev}} - d_{\text{now}}}{d_{\text{prev}}}
        - \beta_{\text{pos}}\, d_{\text{now}}
        - \beta_{\text{rot}}\, y_{\text{now}},
    \qquad \alpha = 3.0,\; \beta_{\text{pos}} = 0.5,\; \beta_{\text{rot}} = 0.25.
    \label{eq:fractional}
\end{equation}

This reward performs poorly for two related reasons. Both concern the gradient it
supplies to PPO, not the optimal policy it defines. Consider a representative push
that moves the object $1\,\text{cm}$ closer to a goal $30\,\text{cm}$ away. The
improvement term contributes $3.0 \times 0.01/0.30 \approx 0.10$. The two penalty
terms contribute $-0.5 \times 0.29 \approx -0.145$ and $-0.25 \times 0.5 =
-0.125$, for a net reward of about $-0.17$. A push that makes genuine progress
therefore receives negative reward, because the standing penalty on the remaining
distance dominates the small improvement. This is the first problem: the reward is
negative for most useful pushes, so the policy gradient points away from
goal-directed behaviour.

The second problem is near-zero variance. Across a rollout, roughly 99\,\% of
pushes produce a reward in the narrow interval $[-0.5, +0.1]$, and the sparse
completion bonuses fire on about 1\,\% of pushes. At 528 environments and 15
pushes per rollout, only about 79 bonus events occur per iteration, too few to
overcome the noise in the dense term. The Generalized Advantage Estimation (GAE)
advantage is then dominated by noise, and the gradient is uninformative. The
failure is therefore one of gradient quality and reward variance, not of a
mis-specified optimum. The reward defines a reasonable objective but supplies a
poor learning signal.

\subsection{Potential-Based Reward Shaping for Pushing}
\label{sec:pbrs_design}

The remedy is to replace the improvement reward with a potential-based reward.
The per-push signal is then dense, correctly signed, and bounded, while the
optimal policy is provably unchanged (Section~\ref{sec:pbrs}). The design uses two
bounded exponential potentials, one for position and one for orientation:
\begin{align}
    \Phi_{\text{pos}}(s) &= w_{\text{pos}}\, \exp\!\left(-k_p\, d^2\right), \label{eq:phi_pos}\\
    \Phi_{\text{rot}}(s) &= w_{\text{rot}}\, \exp\!\left(-k_r\, c(s)\right),
    \qquad c(s) = \frac{1 - \cos(\theta^{*} - \theta)}{2}, \label{eq:phi_rot}
\end{align}
where $d$ is the planar position error and $c(s)$ is a cosine orientation
distance. The cosine form is smooth and free of the discontinuity a modular angle
difference would introduce at $\pm\pi$; $\theta^{*}$ and $\theta$ are the goal and
object yaw. The parameters are $k_p = 30$, $k_r = 5$, and $w_{\text{pos}} =
w_{\text{rot}} = 10$. The dense reward on a transition is the difference of the
summed potential,
\begin{equation}
    F(s, s') = \big[\Phi_{\text{pos}}(s') - \Phi_{\text{pos}}(s)\big]
             + \big[\Phi_{\text{rot}}(s') - \Phi_{\text{rot}}(s)\big],
    \label{eq:pbrs_dense}
\end{equation}
which is Equation~\ref{eq:pbrs} with $\gamma_{\text{shaping}} = 1$. The
undiscounted setting is admissible because the task is episodic and every episode
terminates, as established for episodic MDPs by \cite{Grzes2017RewardShaping}. Two
sparse bonuses are layered on the dense term. A completion bonus of $+5$ is given
when the position gate is satisfied, and a further $+2$ when the combined gate is
satisfied. A penalty of $-5$ is applied if the block tips over. These bonuses are
not themselves potential-based, but they fire only at terminal or near-terminal
events, where they do not create the reward-collecting cycles that
Section~\ref{sec:pbrs} identifies as the danger of non-potential shaping.

On the same representative push, the potential-based reward inverts the sign of
the baseline. For a $1\,\text{cm}$ improvement from $30\,\text{cm}$, the position
potential difference is $10 \times [\exp(-30 \times 0.29^2) - \exp(-30 \times
0.30^2)] \approx +0.13$. Every push toward the goal now receives positive reward,
and every push away from it negative reward. Because the potential is bounded in
$[0, w]$, the reward magnitude cannot be inflated by a physics contact spike.
Figure~\ref{fig:pbrs_signal} contrasts the per-push reward of the two schemes over
a representative episode, and shows the potential landscape and the gradient
comparison.

\begin{figure}[H]
    \centering
    \includegraphics[width=0.48\textwidth]{a3_episode_simulation.png}\hfill
    \includegraphics[width=0.48\textwidth]{a2_gradient_comparison.png}
    \caption{Left: per-push reward over a representative five-push episode for the
    improvement reward and the potential-based reward. The improvement reward is
    negative or near-zero for most pushes; the potential-based reward is
    correctly signed on every push. Right: the gradient supplied by the two
    schemes as a function of distance to the goal.}
    \label{fig:pbrs_signal}
\end{figure}

The reward is implemented as a potential-based reward function. Four properties
make it suitable for a contact-rich task, and all follow from its form. It is
Markovian, since it depends only on the current and successor state. It is
bounded. It produces zero net reward over any cycle that returns to the same
state, which removes the incentive for reward-collecting loops. And it produces
zero reward while the object is stationary, so it imposes no penalty for
deliberation.

\subsection{The SE(2) Unified Distance}
\label{sec:se2}

The two-potential reward treats position and orientation as separate terms whose
gradients can point in different directions. For the self-play models
(Section~\ref{sec:models}) the reward instead uses a single $\mathrm{SE}(2)$
distance, so one potential governs the whole pose. Following the Lie-group
formulation of Section~\ref{sec:pushing}, the pose error is summarised by the
scalar
\begin{equation}
    d_{\text{pose}} = \sqrt{\Delta x^2 + \Delta y^2 + L^2\, \Delta\theta^2},
    \label{eq:dpose_methods}
\end{equation}
where $\Delta x$, $\Delta y$ are the position error components and $\Delta\theta$
the orientation error. Here $L$ is a characteristic length that converts an
angular error into an equivalent linear displacement. For the T-block $L =
0.07\,\text{m}$, the block's radius of gyration; it is a physical property of the
object, not a tuning parameter. The corresponding single potential and its dense
reward are
\begin{equation}
    \Phi_{\text{pose}}(s) = w\, \exp\!\left(-k\, d_{\text{pose}}^2\right),
    \qquad F(s,s') = \Phi_{\text{pose}}(s') - \Phi_{\text{pose}}(s),
    \label{eq:phi_dpose}
\end{equation}
with $k = 30$ and $w = 10$. The success threshold for this metric is
$d_{\text{pose}} < 0.055\,\text{m}$ for the T-block. The single potential removes
the competing-gradient problem of the two-potential form. A push that reduces
$d_{\text{pose}}$ is rewarded whether the reduction comes from position or
orientation, so the reward is aligned with the coupled physics of the task. Both
the two-potential and the unified formulations share the same reward
implementation.

\subsection{The Five Models}
\label{sec:models}

All five models share the same backbone: a single PPO agent with an LSTM of
hidden size 256 followed by a multilayer perceptron, and a MultiCategorical head
over the four-dimensional, 21-bin push primitive. They differ only in the reward
and the training dynamics, which is what makes the comparison controlled. The
five models are summarised in Table~\ref{tab:models}.

\begin{table}[H]
    \centering
    \small
    \begin{tabularx}{\textwidth}{l l X}
        \hline
        \textbf{Model} & \textbf{Training dynamics} & \textbf{Reward} \\
        \hline
        PPO-Baseline & single agent & improvement reward (Equation~\ref{eq:fractional}) \\
        PPO-PBRS & single agent & two-potential PBRS (position, orientation) \\
        PPO-Curriculum & single agent, staged & two-potential PBRS with forced position-then-orientation staging \\
        ASP-dPose & asymmetric self-play & $\mathrm{SE}(2)$ PBRS for Bob, outcome-based Alice reward \\
        TASP-dPose & asymmetric self-play & $\mathrm{SE}(2)$ PBRS for Bob, time-based Alice reward \\
        \hline
    \end{tabularx}
    \caption{The five models. Every model uses the same PPO+LSTM backbone and the
    same object-relative push primitive; only the reward and the training
    dynamics change.}
    \label{tab:models}
\end{table}

\paragraph{PPO-Baseline.}
A single PPO agent trained with the improvement reward of
Equation~\ref{eq:fractional}. It serves as the reference point for the reward
comparison of Section~\ref{sec:frac_reward} and is trained at 2{,}048 parallel
environments, a larger batch than the other models use.

\paragraph{PPO-PBRS.}
A single PPO agent trained with the two-potential reward of
Equations~\ref{eq:phi_pos} to \ref{eq:pbrs_dense}. This is the simplest model that
uses the potential-based reward, with no curriculum and no second agent, and it
is trained at 528 environments.

\paragraph{PPO-Curriculum.}
PPO-PBRS with a forced three-phase curriculum that stages the two objectives.
Phase 1 trains position only, with the orientation weight set to zero. Phase 2
begins when the episodic position-only success rate reaches 0.5 over a
20-iteration window; the orientation weight is then increased from 0 to 10 over
200 iterations. Phase 3 is the full two-objective reward. The trigger gates on the
episodic position-only success rate, rather than on a per-push error mean. This is
the corrected version of a trigger that in an earlier implementation never
activated; the earlier failure and its correction are documented in Appendix~H.

\paragraph{ASP-dPose.}
An asymmetric self-play model that uses the $\mathrm{SE}(2)$ reward of
Equation~\ref{eq:phi_dpose} for the goal-solving agent, Bob. The loop is shown in
Figure~\ref{fig:asp_loop}. The goal-setting agent, Alice, performs five pushes to
propose a goal. The goal is accepted only if it passes a validity check: the block
has moved sufficiently, remains on the table, and lies inside the workspace. Bob
then attempts the goal from a clean reset over up to ten pushes. Alice's reward is
\emph{outcome-based}, following \cite{OpenAI2021AsymmetricSF}: she receives a
reward when Bob fails, a shaping bonus for a valid goal, and a penalty for an
out-of-bounds goal. Bob additionally uses a GoalEncoder, a small multilayer
perceptron that maps the six-dimensional goal pose into an eight-dimensional
latent, in the difference variant of \cite{Sukhbaatar2018GoalEmbeddings}. Bob is
trained with Alice Behavioral Cloning at $\beta = 0.5$.

\paragraph{TASP-dPose.}
Identical to ASP-dPose except that Alice's reward is \emph{time-based}, following
\cite{Sukhbaatar2018IntrinsicMotivation}: $R_A = \gamma_{sp}\max(0, t_B - t_A)$.
This rewards Alice for proposing goals that take Bob more pushes to solve than
Alice took to set, subject to a penalty on Alice's own effort. It is the
self-regulating reward of Equation~\ref{eq:sukh_alice}. The two self-play models
therefore differ only in Alice's reward, which isolates the effect of the goal
proposal mechanism.

\begin{figure}[H]
    \centering
    \includegraphics[width=0.9\textwidth,height=0.5\textheight,keepaspectratio]{asp_loop_diagram.png}
    \caption{The asymmetric self-play loop. Alice proposes a goal by acting for a
    fixed number of pushes; the goal is validated; Bob then attempts it from a
    clean reset. Alice's reward differs between ASP-dPose (outcome-based) and
    TASP-dPose (time-based); Bob's $\mathrm{SE}(2)$ potential-based reward is the
    same in both.}
    \label{fig:asp_loop}
\end{figure}

\subsection{Why PBRS Assists Self-Play but Does Not Solve It}
\label{sec:pbrs_asp}

Potential-based shaping improves self-play but does not make it competitive with
single-agent training, and the reason is visible in the structure of the GAE
advantage. The advantage estimate is a discounted sum of one-step temporal
differences, each of which decomposes into a reward term and a value-bootstrap
term:
\begin{equation}
    \hat{A}_t = \sum_{l \geq 0} (\gamma\lambda)^l
    \underbrace{\big[\Phi(s_{t+l+1}) - \Phi(s_{t+l})\big]}_{\text{potential term, correctly signed}}
    \;+\; \sum_{l \geq 0} (\gamma\lambda)^l
    \underbrace{\big[\gamma V(s_{t+l+1}) - V(s_{t+l})\big]}_{\text{value term, biased when } V \text{ is stale}}.
    \label{eq:gae_decomp}
\end{equation}
The potential term is correctly signed on every push, regardless of the accuracy
of the value function, because the potential increases whenever the object moves
toward the goal. The value term, by contrast, depends on the accuracy of $V$. In
single-agent training the goal distribution is stationary, so $V$ is accurate and
both terms agree. Under self-play, Alice changes the goal distribution at every
iteration, so $V$ is perpetually estimated on a distribution that no longer holds,
and the value term is biased. The potential term supplies a floor of correct
learning signal, which lifts self-play above the level it reaches under the sparse
or improvement reward. It cannot, however, remove the bias contributed by the
stale value term. The consequence, quantified in Chapter~\ref{sec:results}, is
that
potential-based shaping raises self-play from near-zero to a modest success rate
without approaching the single-agent result.

This account is a hypothesis about the learning dynamics, not a result. Both of
its claims, the elevated value error under a shifting goal distribution and the
suppression of the imitation signal, are measured directly through the training
instrumentation described in Section~\ref{sec:instrumentation}, and the
measurements are reported in Chapter~\ref{sec:results} rather than asserted here.

\subsection{Experimental Campaign Overview}
\label{sec:campaign}

The five models of Section~\ref{sec:models} are the objects of study. A single
training run of each, at one batch size and one seed, cannot support the
statistical and causal claims this thesis makes. The evaluation is therefore
organised as a campaign of budget-matched training runs, summarised in
Table~\ref{tab:campaign}. Its per-phase configurations and job counts are given in
Chapter~\ref{sec:experiments}. This section states the two principles that make
the campaign a controlled study, together with the validation protocol and
training instrumentation that every run shares. The numerical outcomes appear in
Chapter~\ref{sec:results}.

The first principle is \textbf{multiple seeds with reported confidence intervals}.
Single-seed reinforcement-learning results are known to be unreliable, because the
variance between random seeds can exceed the difference between methods
\cite{Henderson2018DeepRLMatters, Agarwal2021StatisticalPrecipice}. Every headline
comparison in this thesis is therefore repeated across several seeds. Each is
reported as a mean with a 95\,\% confidence interval, so a claimed difference
between two models is distinguished from seed noise.

The second principle is \textbf{budget matching}. The aim is to compare batch
sizes without confounding them with the amount of experience. The total number of
pushes seen over training is therefore held constant while the number of parallel
environments is varied. The iteration count is set inversely proportional to the
number of environments, so each configuration consumes the same total experience
of about 23.8 million pushes. The temporal window of the recurrent policy is fixed
at 15 pushes per rollout; only the number of parallel environments, and therefore
the per-update batch size, changes. This isolates the per-update batch size, the
quantity that drives the value-function bias of Section~\ref{sec:pbrs_asp}, from
the total experience.

\begin{table}[H]
    \centering
    \small
    \begin{tabularx}{\textwidth}{l l X}
        \hline
        \textbf{Phase} & \textbf{Question} & \textbf{Design} \\
        \hline
        Scale sweep & Does the self-play deficit close as the batch grows? &
        the two self-play models across four environment counts \\
        Environment frontier & How few environments does each single-agent reward need? &
        potential-based single agent and the improvement-reward baseline across four environment counts, three seeds \\
        Seed intervals & Are the model differences larger than seed noise? &
        single agent, curriculum, and both self-play models at the anchor budget, five seeds \\
        Self-play ablations & Which component of the self-play package matters? &
        the self-play model with the goal encoder, behavioural cloning, or the historical pool removed one at a time, three seeds \\
        Positive control & Does behavioural cloning help on an easy symmetric task? &
        self-play on the rotationally symmetric disc with behavioural cloning enabled, three seeds \\
        Reward sensitivity & Is the potential-based reward robust to its parameters? &
        the single agent over a grid of the position sharpness $k_p$ and the potential weight $w$ \\
        Cross-environment & Does the ordering hold in a second simulator? &
        single agent, curriculum, and self-play in the two-dimensional pushing simulator across two batch sizes \\
        \hline
    \end{tabularx}
    \caption{The experimental campaign. Each phase answers one question flagged in
    review; the per-phase training configurations, environment counts, seeds, and
    job counts are tabulated in Chapter~\ref{sec:experiments}.}
    \label{tab:campaign}
\end{table}

\subsection{Validation Protocol}
\label{sec:val_protocol_methods}

Every trained policy is evaluated on the same held-out suite of 30 T-block
scenes, illustrated in Figure~\ref{fig:val_layout_methods}. The suite is divided
into three groups of ten. The first ten are rotation-dominant scenes with large
orientation changes. The next ten are position-only scenes whose goal orientation
equals the start orientation. The last ten are combined scenes that require both a
position and an orientation change. Within each group the difficulty ranges from
easy to hard, graded by the start-to-goal distance and the required rotation. The
held-out scenes are fixed in advance and never seen during training, so
performance on them measures generalisation rather than memorisation.

\begin{figure}[H]
    \centering
    \includegraphics[width=\textwidth,height=0.42\textheight,keepaspectratio]{val_test_layout.png}
    \caption{The 30-scene held-out validation suite. Each panel shows the block
    start pose and the goal pose for one scene. The suite is identical for every
    model and every seed.}
    \label{fig:val_layout_methods}
\end{figure}

Each scene is attempted 20 times, and within each attempt the policy is allowed
up to 30 pushes. Two success measures are reported and kept distinct throughout
the thesis. The \textbf{scene success rate} counts a scene as solved if at least
one of its 20 attempts satisfies the combined gate of Section~\ref{sec:mdp}. This
is the headline measure, and follows the convention common in manipulation
evaluation, where a task counts as achievable if the policy can achieve it at all.
The \textbf{trial success rate} is the fraction of individual attempts that satisfy
the gate, and is always the lower of the two. Reporting both prevents the
best-of-20 protocol from overstating performance, and the gap between them is
stated wherever the headline number appears.

Three protocol decisions make the numbers comparable across models. Each corrects
a subtle way the comparison could otherwise be invalidated. First, the checkpoint
selected for validation is the one that maximised the \emph{combined}
position-and-orientation success rate during training, not the looser
single-distance threshold used internally by the self-play reward. Selecting on a
different criterion than the one reported would flatter the self-play models.
Second, the self-play policies trained on the rotationally symmetric disc are
validated on disc scenes, whose goal orientation is unconstrained, not on the
T-block scenes they never trained on. Validating an agent on a task it never
trained for would void its result. Third, checkpoints that lack a combined-gate
best model, and are validated only from their latest checkpoint, are recorded
separately and excluded from the confidence-interval aggregation. A
partially-trained run therefore cannot contaminate a headline mean. The per-seed
results are aggregated into a mean with a 95\,\% confidence interval per model,
batch size, and iteration count.

\subsection{Training Instrumentation for the Self-Play Mechanism}
\label{sec:instrumentation}

The account in Section~\ref{sec:pbrs_asp} of why potential-based shaping only
partially assists self-play rests on two claims. The first is that the
value-function term of the advantage estimate is biased under the shifting goal
distribution. The second is that the behavioural-cloning signal is suppressed when
the goal-proposing agent's demonstrated actions lie far from the goal-solving
agent's policy. Rather than leave these as arguments, the training runs are
instrumented so that both are measured directly, and the measurements are reported
in Chapter~\ref{sec:results}.

For the value-function claim, the mean absolute difference between the value
estimates and the empirical returns is logged at every policy update. This is done
for the single-agent, self-play, and time-based self-play models alike. Under the
budget-matched scale sweep, this quantity can be plotted against the number of
parallel environments. The mechanism of Section~\ref{sec:pbrs_asp} makes two
predictions: the value error is larger under self-play than under single-agent
training at the same batch size, and it falls as the batch grows. A value error
that declines toward the single-agent level with scale would support a
budget-limitation account. One that remains elevated would indicate a structural
deficit independent of budget.

For the behavioural-cloning claim, the imitation term is instrumented in three
ways: the share of demonstrated actions whose probability ratio falls outside the
clipping interval, the mean ratio, and the mean probability the goal-solving agent
assigns to the goal-proposing agent's demonstrated actions. The clipping mechanism
inherited from the imitation loss of \cite{OpenAI2021AsymmetricSF} zeroes the
gradient for any demonstrated action whose ratio leaves the trust interval. A large
clipped fraction is therefore the direct signature of a suppressed imitation
signal. Measuring this fraction converts the clip-saturation account into a
reported quantity. A positive control on the rotationally symmetric disc, where
behavioural cloning is enabled and the task is easy, tests whether the mechanism
can help at all in this setting. If imitation assists anywhere it should assist
there, so its failure on the coupled T-block task cannot then be dismissed as a
broken implementation.
